\section{Introduction}\label{sec:intro}

% TODO (author): add 2--3 current epidemiology references for this opening claim (paper_structure.md §1).
Alzheimer's disease (AD) and frontotemporal dementia (FTD) are leading causes of neurodegenerative
dementia, particularly with onset before 65, and telling them apart is clinically consequential:
they differ in prognosis, in care needs and in which disease-modifying treatments are appropriate.
Scalp electroencephalography (EEG) is attractive for this task because it is cheap, non-invasive and
available in routine care, in contrast to amyloid PET or MRI. Its physiological basis is well
established: AD is increasingly understood as a disorder of large-scale brain networks
\cite{F12}, and graph-theoretical EEG studies have shown that AD and FTD disturb those networks in
\emph{different} ways \cite{F9}, which is what makes a three-class AD/FTD/control problem
scientifically meaningful rather than merely harder.

Machine-learning work on this question has moved from spectral markers and handcrafted features with
classical classifiers \cite{F10} to convolutional networks and transformers \cite{F2,R14} and,
most recently, to graph neural networks (GNNs) that treat electrodes as nodes and inter-electrode
coupling as edges \cite{F5,F3,F4,F6,R9,R10,R12}. The public OpenNeuro dataset ds004504
\cite{F1,ds004504}, with 88 subjects recorded on one device under one protocol, has become the
de-facto benchmark. Reported accuracies on this single dataset span roughly 80--99\%.

That spread is implausibly wide for identical data, and it hides a methodological problem. Published
studies change several things at once: (i) the \emph{signal representation} fed to the model,
(ii) the \emph{graph construction} that defines which electrodes are connected, (iii) the
\emph{model architecture}, and (iv) the \emph{evaluation protocol}. Independent evidence suggests
that the last of these matters most. A 2026 scoping review of this exact benchmark concluded that
apparent deep-learning gains ``arise not from inherent representational advantage, but from data
leakage'' \cite{R1}, echoing a broader warning about leakage across machine-learning-based science
\cite{F8} and in neuroimaging prediction specifically \cite{L5}. In clinical machine learning,
splitting records rather than subjects has long been known to overestimate accuracy for new
patients \cite{L1}, and segment-level splitting is common in, and inflates, deep-learning studies of
translational EEG \cite{L2,L3}. A 2026 comparison of five time--frequency representations under a single convolutional
network \cite{R2} isolated representation, but its authors note that it ``does not explicitly model
the true spatial topology'' and that graph-based schemes remain unexplored; others observe that most
studies ``overlook spatial connectivity between channels'' \cite{R9}, and explicitly call for graph
convolutional networks in future work \cite{R8,F7}. No study has varied representation and graph
construction together while holding the model and the evaluation protocol fixed.

A second, biologically motivated gap concerns the blood--brain barrier (BBB). BBB breakdown is an
early biomarker of cognitive dysfunction \cite{B5,B8}, and Milikovsky et al.\ described
\emph{paroxysmal slow-wave events} (PSWEs) --- episodes in which the EEG median power frequency
falls below \SI{6}{\hertz} for at least \SI{5}{\second} --- in patients with AD, linked them to
cognitive impairment, localised them to cortex with BBB dysfunction in epilepsy, and induced them in
rodents by albumin extravasation \cite{B1,B2}. PSWEs have since been examined in Parkinson's disease
\cite{B3} and epilepsy \cite{B4}, but, to our knowledge, never in FTD and never on an open
dementia cohort. Whether this BBB-associated EEG marker carries diagnostic information beyond
ordinary EEG slowing has not been tested.

We therefore ran a controlled factorial experiment. Seven node-feature representations (spectral,
connectivity, wavelet, time-domain, raw, fused, and PSWE-derived) were crossed with three graph
constructions (spatial, functional, hybrid) under a single fixed two-layer graph convolutional
network \cite{F13}, with subject-level repeated cross-validation, identical frozen folds for every
condition, and hyperparameters fixed before any result was seen. Because the model is the control
variable rather than the contribution, it is deliberately small. Around this core we measured the
inflation produced by window-level (leaky) splitting on the identical pipeline, quantified
age/sex confounding, and characterised PSWE burden across AD, FTD and controls.

Our contributions are:
\begin{enumerate}
  \item \textbf{Methodological.} A fully crossed $7\times3$ protocol that separates signal
        representation, graph topology and model architecture, with paired statistics over 25
        frozen subject-level test splits and ablations that test whether the graph matters at all.
  \item \textbf{Empirical.} Effect sizes and confidence intervals for which representations a GNN
        benefits from in three-class AD/FTD/control classification, with non-graph twins and a
        demographic confound floor as reference points.
  \item \textbf{Protocol.} A direct measurement, on one pipeline, of how much window-level splitting
        inflates results --- a calibration factor for reading published numbers on this benchmark.
  \item \textbf{Biological.} The first characterisation of BBB-associated PSWEs in FTD, an
        independent-cohort test of the AD finding, and a test of whether the marker is separable
        from background slowing.
  \item \textbf{Reproducibility.} Code, seeds, frozen folds and every raw run released; the full
        study re-runs on free cloud CPUs.
\end{enumerate}
